\documentclass[10pt,a4paper]{amsart}
\usepackage[margin=19mm]{geometry}
\usepackage{amsmath,amssymb,mathtools}
\usepackage[scr=rsfs,cal=euler]{mathalfa}
\usepackage{xfrac}
\usepackage[unicode,hidelinks,bookmarks=false]{hyperref}
\newcommand{\ie}{i.e.\ }
\newcommand{\cf}{cf.\ }
\newcommand{\resp}{respectively\ }
\newcommand{\Subsection}{\subsection}
\providecommand{\nobreakdash}{\nobreak\hbox{-}}
\setlength{\emergencystretch}{2em}
\multlinegap0pt
\usepackage{amssymb}
\usepackage{xcolor}
\usepackage{enumerate}
\usepackage{tikz}
\newtheorem{theorem}{Theorem}
\title{On the revolving structure of the L\'evy dragon and its linear transforms}
\author{Miguel Gonzalez-Carriedo, Kiko Kawamura, Jonathan Leung, R. D\'aniel Prokaj}
\date{Source: arXiv:2606.05565v1 (CC BY 4.0)}
\begin{document}
\maketitle

\noindent\emph{Source and licence.} On the revolving structure of the L\'evy dragon and its linear transforms. Version arXiv:2606.05565v1 (CC BY 4.0), distributed by arXiv under the Creative Commons Attribution 4.0 International licence, http://creativecommons.org/licenses/by/4.0/, as recorded in the arXiv licence metadata for this exact version and shown on the abstract page https://arxiv.org/abs/2606.05565v1. Full licence notice: \texttt{licenses/arxiv-2606.05565-CC-BY-4.0.txt}.\par\smallskip

\noindent\textit{Editorial scope.} Counted unit: the article's theorem lt93 (source0.tex lines 263--267). The article's complete original proof of it is delivered with it (source0.tex lines 379--418, one \texttt{proof} environment of 40 lines, which the article places away from the statement). No mathematics is written, repaired or re-derived: the statement and the proof are the article's own lines, in the article's own notation, and no retained line is substituted or re-flowed.  Retained uncounted as the source-local context the proof needs: lines 300--303.  One declared adaptation: exactly 1 ancillary strings inside the retained spans are omitted (image-inclusion commands and, where the source repeats a label, the repeated label definitions) and no other line differs from the source; the figure environments, with their placement, captions and labels, are kept, and the package records the omitted strings in fidelity receipts bound to the affected bodies.  No retained line contains a citation, so the wrapper carries no bibliography and no bibliography entry.  The retained lines contain the article's own diagram code, which the wrapper typesets with the standard tikz packages; no image file is delivered and the package declares no asset dependency.  Editorial formatting only, in addition: an AMS \texttt{amsart} preamble that restates the article's own declarations and macros used by the retained lines, namely the environments \texttt{theorem} on separate counters, together with the standard packages those lines need.  The retained environments print in this document as Theorem 1 in order of appearance, whereas the article prints the same items with the numbering of its own full document.  The complete native source of the article as distributed in the arXiv e-print for this exact version is bundled unchanged as \texttt{sources/202126/source0.tex}.
\begin{equation}\label{lt98}
    \Pi(\pmb{\omega})=\left(\frac{1-i}{2}\right)^k i^{q_k(\pmb{\omega})} \Pi(\sigma(\pmb{\omega})) +
    \sum_{n=1}^{k}\omega_n\left(\frac{1-i}{2}\right)^n i^{q_{n-1}(\pmb{\omega})}.
\end{equation}

\begin{theorem}[Main Theorem]\label{lt93}
  Let $S$ be a non-singular linear transformation on $\mathbb{C}$ and $S(\Lambda)$ be the transformed L\'evy dragon. Then, $S(\Lambda)$ also admits a revolving structure.
  Further, $\Lambda$ and $S(\Lambda)$ have the same revolving pattern.

\end{theorem}

\begin{proof}[Proof of Theorem~\ref{lt93}]
    Let $S(z)=\lambda z+\tau$ be a non-singular linear transformation on $\mathbb{C}$. Then $S(\Lambda)$ is the attractor of the IFS 
    \begin{equation}
      \widehat{\mathcal{F}}=\left\{
        \widehat{f}_0(z)= S\circ f_0\circ S^{-1}(z), \quad
        \widehat{f}_1(z)= S\circ f_1\circ S^{-1}(z)
      \right\}.
    \end{equation}
    Using that $S^{-1}(z)=\frac{z-\tau}{\lambda}$, a simple calculation gives
    \begin{equation}
      \widehat{\mathcal{F}}=\left\{
        \widehat{f}_0(z)= \frac{1-i}{2}z+\frac{1+i}{2}\tau, \quad
        \widehat{f}_1(z)= \frac{1+i}{2}z+\frac{1-i}{2}(\lambda+\tau) 
      \right\}.
    \end{equation}

    Analogously to \eqref{lt98}, the canonical projection $\widehat{\Pi}$ of $\widehat{\mathcal{F}}$ can be written in the form
    \begin{align*}
      \widehat{\Pi}(\pmb{\omega}) &=\left(\frac{1-i}{2}\right)^k i^{q_k(\pmb{\omega})}\widehat{\Pi}(\sigma^k\pmb{\omega})\\ 
      &+ \sum_{n=1}^{k} \left(\frac{1-i}{2}\right)^{n-1} i^{q_{n-1}(\pmb{\omega})} \left(\frac{1+i}{2}\tau +\omega_n\left(\frac{1-i}{2}\lambda -i\tau\right)\right).
    \end{align*}
    Then, we take the limit on both sides as $k\to\infty$ to obtain
    \[
      \widehat{\Pi}(\pmb{\omega})=\sum_{n=1}^{\infty} \left(1+i\right)^{-n} i^{q_{n-1}(\pmb{\omega})} \left( i\tau + \omega_n\left(\lambda +(1-i)\tau\right)\right).
    \]
    It follows that the points of the transformed L\'evy dragon $S(\Lambda)$ are the revolving numbers of base $1+i$ and digit set 
    \[
      \left\{ 
        \lambda+\tau,\: i(\lambda+\tau),\: -(\lambda+\tau),\: -i(\lambda+\tau),\: 
        \tau,\: i\tau,\: -\tau,\: -i\tau
      \right\}.
    \]
    Further, the digits must follow each other according to the directed graph on Figure \ref{lt94}. Note that the first digit $\delta_1$ is either $\lambda+\tau$ or $i\tau$. 

    \begin{figure}
    
    \caption{The directed graph of the revolving structure of $S(\Lambda)$.}\label{lt94}
    \end{figure}

\end{proof}

\end{document}
